\documentclass[12pt]{article}
\usepackage[utf8]{inputenc}
\usepackage[T1]{fontenc}
\usepackage{hyperref}
\title{Chepya Language Resources for K-12 Curriculum Development}
\author{Compiled by Claude Sonnet 5.5 (claude-sonnet-5-5) for Andr\'as Kornai}
\date{}
\begin{document}
\maketitle

\section{Introduction}
Chepya (ISO 639-3 ycp; Glottocode chep1243) is a Loloish (Sino-Tibetan) language of northern Laos, centered on Phongsaly Province. Wikipedia gives about 2,000 speakers as of 2007. Joshua Project lists a population of 2,700 for the Chepya people. Very little exists for this language. The search found only general reference pages.

\section{Primary Educational Resources}
None were found. Chepya is not known to be taught in schools.

\section{Linguistic and Reference Materials}
\begin{itemize}
\item Wikipedia, Chepya language. It cites a 1999 vocabulary collection by Kingsada and Shintani as the main description: \url{https://en.wikipedia.org/wiki/Chepya_language}
\item Glottolog entry: \url{https://glottolog.org/resource/languoid/id/chep1243.bigmap.html}
\item Verbix endangered-languages page, which classifies Chepya as threatened: \url{https://docs.verbix.com/EndangeredLanguages/Chepya}
\end{itemize}

\section{Religious and Community Resources}
\begin{itemize}
\item Joshua Project people-group and language pages. Joshua Project describes the language as Mon-Khmer, which conflicts with the linguistic sources, so treat that detail with caution: \url{https://joshuaproject.net/people_groups/20619} and \url{https://joshuaproject.net/languages/ycp}
\end{itemize}
No scripture or community organisation was found.

\section{Digital and Audio Resources}
None were found.

\section{Curriculum Development Context}
There is essentially no existing educational infrastructure for Chepya. Any curriculum work would start with documentation and community consultation.
\end{document}